\begin{pastexam}[label={px:ntut-115-ee-ds-22}]{北科大電機所（戊組）115 學年度 問答題第 2 題}
（問答）

Consider the following undirected graph:\\
Vertices: A, B, C, D, E, F\\
Edges: (A,B), (A,C), (B,D), (C,D), (C,E), (D,E), (D,F), (E,F)\\
(a) Draw the adjacency matrix representation of this graph. \textbf{(5\%)}\\
(b) Starting from vertex A, show the order of vertices visited using Breadth-First Search (BFS). Assume
that when multiple vertices are available, they are visited in alphabetical order. \textbf{(5\%)}
\end{pastexam}

\begin{pastsoln}
\textbf{(a)} 鄰接矩陣（資結書第~21.2~節）：有邊 $(u, v)$ 時 $M[u][v] = M[v][u] = 1$，否則 $0$。 無向圖的矩陣對稱，
對角線是 $0$（沒有自迴圈）。
\begin{center}
\begin{tabular}{c|cccccc}
 & A & B & C & D & E & F \\
\hline
A & 0 & 1 & 1 & 0 & 0 & 0 \\
B & 1 & 0 & 0 & 1 & 0 & 0 \\
C & 1 & 0 & 0 & 1 & 1 & 0 \\
D & 0 & 1 & 1 & 0 & 1 & 1 \\
E & 0 & 0 & 1 & 1 & 0 & 1 \\
F & 0 & 0 & 0 & 1 & 1 & 0 \\
\end{tabular}
\end{center}
檢查：共 $16$ 個 $1$，是邊數 $8$ 的兩倍；每列的和是度數，A 到 F 依序是 $2, 2, 3, 4, 3, 2$，加起來 $16 = 2 \times 8$。

\textbf{(b)} BFS 的拜訪順序：\textbf{A, B, C, D, E, F}。

用佇列（資結書第~21.3~節）：取出一個頂點，把它還沒拜訪過的鄰居依字母順序拜訪並放進佇列。
\begin{center}
\small
\begin{tabular}{llll}
\toprule
取出 & 鄰居（字母順序） & 新拜訪的 & 佇列（前 $\to$ 後） \\
\midrule
（開始） & & A & A \\
A & B, C & B, C & B, C \\
B & A, D & D & C, D \\
C & A, D, E & E & D, E \\
D & B, C, E, F & F & E, F \\
E & C, D, F & — & F \\
F & D, E & — & （空） \\
\bottomrule
\end{tabular}
\end{center}
各頂點離 A 的距離：A 是 $0$；B、C 是 $1$；D、E 是 $2$；F 是 $3$。 BFS 依距離一層一層拜訪，同一層內照字母順序，
所以順序是 A、B、C、D、E、F。 BFS 樹的邊是 (A,B)、(A,C)、(B,D)、(C,E)、(D,F)。
\end{pastsoln}
